% ============================================================================
% บทที่ 3 — ตั้งค่าและนำเข้าข้อมูล
% ============================================================================
\mychapter{ตั้งค่าและนำเข้าข้อมูล}

\noindent บทนี้อธิบายการเพิ่ม/ลบไฟล์ การกำหนดบทบาท ตั้งค่าปฏิทิน และการนำเข้าข้อมูล

\section{การเพิ่มและลบไฟล์รายวิชา}

\noindent\textbf{การเพิ่มไฟล์}
\begin{steps}
  \item ไปที่หน้า \textbf{ข้อมูลและการตั้งค่า}
  \item ในส่วน \textbf{ไฟล์รายวิชา} ให้กดปุ่ม \texttt{เพิ่มไฟล์รายวิชา}
  \item เลือกไฟล์ข้อมูลที่ต้องการ (สามารถเลือกได้หลายไฟล์พร้อมกัน)
\end{steps}

\noindent\textbf{การลบหรือเปลี่ยนไฟล์}
\begin{steps}
  \item หากต้องการลบไฟล์ ให้กดปุ่มรูปถังขยะ \texttt{ลบไฟล์} ที่อยู่ขวาของชื่อไฟล์
  \item หากต้องการเปลี่ยนไฟล์ ให้ลบไฟล์เดิมแล้วเพิ่มไฟล์ใหม่
\end{steps}

\section{การกำหนดบทบาทของไฟล์}

สำหรับแต่ละไฟล์ในส่วนไฟล์รายวิชา ให้เลือกบทบาทจากช่องแบบเลื่อนลงถัดจาก
ชื่อไฟล์ (\texttt{จัดอัตโนมัติ} หรือ \texttt{จัดไปแล้ว}) โดยค่าเริ่มต้นของทุกไฟล์
คือ \texttt{จัดอัตโนมัติ}

\section{การเพิ่มไฟล์เงื่อนไข}

ไฟล์เงื่อนไขการสอบเป็นตัวเลือกที่ไม่บังคับ:

\begin{steps}
  \item ในส่วน \textbf{เงื่อนไขการสอบ} ให้กดปุ่ม \texttt{เลือกไฟล์}
  \item เลือกไฟล์ \file{rules.csv} หรือไฟล์เงื่อนไขที่ต้องการ
\end{steps}

\section{การตั้งค่าปฏิทิน}

ในส่วนปฏิทินการสอบ ให้กำหนดช่วงเวลาสอบของทั้งสองภาค โดยระบุ
\textbf{วันเริ่มสอบกลางภาค} และ \textbf{วันสิ้นสุดสอบกลางภาค} เช่น
\texttt{2026-08-17} ถึง \texttt{2026-08-28} และทำเช่นเดียวกันกับ
\textbf{วันเริ่มสอบปลายภาค} และ \textbf{วันสิ้นสุดสอบปลายภาค} เช่น
\texttt{2026-10-14} ถึง \texttt{2026-10-30}

\begin{note}
หากป้อนวันที่เริ่มหลังวันที่สิ้นสุด ระบบจะปฏิเสธ และปุ่มนำเข้าจะไม่ทำงาน
\end{note}

\section{นโยบายวันเสาร์–อาทิตย์และวันหยุด}

ระบบมีนโยบายสำหรับวันเสาร์–อาทิตย์และวันหยุดสามแบบ:

\begin{description}
  \item[\texttt{ใช้เมื่อจำเป็น} (ค่าเริ่มต้น)] ใช้ได้เมื่อจำเป็นเท่านั้น
  \item[\texttt{ไม่ใช้วันเหล่านี้}] ห้ามใช้โดยเด็ดขาด
  \item[\texttt{ใช้ได้ตามปกติ}] ใช้ได้เหมือนวันธรรมดา
\end{description}

เลือกนโยบายที่ต้องการในช่อง \textbf{วันเสาร์–อาทิตย์} และช่อง \textbf{วันหยุด}
ปฏิทินจะแสดงป้ายกำกับวันเสาร์–อาทิตย์และวันหยุดทุกกรณี แต่แสดงว่าห้ามใช้หรือไม่
ขึ้นอยู่กับนโยบายที่เลือกเท่านั้น

\begin{note}
ระบบไม่อนุมานวันหยุดของมหาวิทยาลัยหรือวันหยุดราชการโดยอัตโนมัติ
ต้องระบุวันหยุดที่ต้องการเองในส่วน \textbf{เพิ่มวันหยุด}
\end{note}

\section{การเพิ่มวันหยุด}

ในส่วนวันหยุด ให้กดปุ่ม \texttt{เพิ่ม} แล้วระบุวันที่ที่เป็นวันหยุด
เช่น \texttt{2026-08-24} และ \texttt{2026-10-21}

\section{การตั้งค่าช่วงเวลาและระดับความพยายาม}

ในหน้าข้อมูลและการตั้งค่า ให้กดเปิดส่วน \textbf{ช่วงเวลาและการค้นหา}
เพื่อดูตัวเลือกขั้นสูง ได้แก่ ช่วงเวลาสอบ ช่วงสอบเต็มวัน และระดับความพยายาม:

\begin{steps}
  \item ตรวจสอบช่วงเวลาสอบที่ใช้ในการจัดอัตโนมัติ:
        \texttt{09:00--12:00} และ \texttt{13:00--16:00}
  \item ตรวจสอบช่วงสอบเต็มวัน: \texttt{09:00--16:00}
  \item ตรวจสอบ \textbf{ระดับความพยายามในการจัดตาราง}
\end{steps}

ช่วงเวลาเหล่านี้เป็นช่วงที่ระบบใช้ในการจัดอัตโนมัติเท่านั้น
เมื่อแก้ไขเวลาด้วยตนเองในบทที่ 5 สามารถป้อนช่วงเวลาอื่นได้

\textbf{ระดับความพยายาม} คือจำนวนรอบที่ระบบพยายามจัดตาราง
แนะนำให้ใช้ค่าเริ่มต้นก่อน หากระบบจัดตารางได้ไม่ครบ อาจเพิ่มค่านี้แล้วลองใหม่
ค่าสูงขึ้นอาจใช้เวลาคำนวณนานขึ้น (ช่วงค่าที่ยอมรับอยู่ในภาคผนวก A)

\section{การนำเข้าและตรวจสอบ}

เมื่อเตรียมไฟล์และตั้งค่าเรียบร้อยแล้ว กดปุ่ม \texttt{นำเข้าและตรวจสอบ}

\manualfigure{docs/usage-report/screenshots/figures/F01.png}{0.95\textwidth}{%
  หน้าข้อมูลและการตั้งค่า แสดงส่วนไฟล์รายวิชา เงื่อนไขการสอบ และปฏิทินการสอบ}

\begin{expected}
ระบบนำเข้าและตรวจสอบข้อมูล หากสำเร็จจะแสดงข้อความ \textbf{นำเข้าเสร็จแล้ว กรุณาตรวจสอบผลก่อนจัดตาราง}
และนำไปยังหน้าตรวจสอบข้อมูล หากข้อมูลมีปัญหา จะแสดงรายการข้อควรตรวจสอบ
\end{expected}

\seealso{บทที่ 6 สำหรับการตรวจสอบและแก้ปัญหาข้อมูล}
